\documentclass[14pt]{extarticle}

% 16:9 page (presentation aspect ratio) - tightened margins
\usepackage[
  paperwidth=257mm,
  paperheight=144.5mm,
  top=9mm,
  bottom=7mm,
  left=5mm,
  right=5mm,
  headheight=7mm,
  headsep=2mm,
  footskip=4mm
]{geometry}

\usepackage{amsmath,amsthm,amssymb,amsfonts}
\usepackage{mathtools}
\usepackage{booktabs,array}
\usepackage[shortlabels]{enumitem}
\usepackage{xcolor}
\usepackage{fancyhdr}
\usepackage[most]{tcolorbox}
\tcbuselibrary{theorems,skins,breakable}
\usepackage{hyperref}
\hypersetup{colorlinks=true,linkcolor=blue!60!black,urlcolor=blue!70!black,citecolor=blue!60!black}
\usepackage{cleveref}

%%% COLOR PALETTE %%%
\definecolor{hdrBG}{HTML}{1E3A8A}        % header blue
\definecolor{accent}{HTML}{0EA5E9}       % subsection cyan
\definecolor{defBG}{HTML}{EFF6FF}        % definition: soft blue
\definecolor{defFR}{HTML}{2563EB}
\definecolor{thmBG}{HTML}{FEF3C7}        % theorem: soft amber
\definecolor{thmFR}{HTML}{B45309}
\definecolor{lemBG}{HTML}{FEF9C3}        % lemma: lighter amber
\definecolor{lemFR}{HTML}{A16207}
\definecolor{corBG}{HTML}{FDF2F8}        % corollary: soft pink
\definecolor{corFR}{HTML}{BE185D}
\definecolor{propBG}{HTML}{F3E8FF}       % proposition: soft purple
\definecolor{propFR}{HTML}{7C3AED}
\definecolor{proofBG}{HTML}{F8FAFC}      % proof: very soft gray
\definecolor{proofFR}{HTML}{64748B}
\definecolor{remBG}{HTML}{ECFDF5}        % remark: soft green
\definecolor{remFR}{HTML}{047857}
\definecolor{notBG}{HTML}{F1F5F9}        % notation: slate
\definecolor{notFR}{HTML}{475569}

%%% HEADER / FOOTER %%%
\pagestyle{fancy}
\fancyhf{}
\renewcommand{\headrulewidth}{0.4pt}
\renewcommand{\headrule}{\color{hdrBG!30}\hrule width\textwidth height0.4pt}
\renewcommand{\footrulewidth}{0pt}
\fancyhead[L]{\small\color{hdrBG}\textbf{Gaussian Correlation Inequality}}
\fancyhead[R]{\small\color{hdrBG}\leftmark}
\fancyfoot[R]{\small\color{gray}\thepage}
\renewcommand{\sectionmark}[1]{\markboth{\S\thesection~#1}{}}

%%% AMSTHM ENVIRONMENTS %%%
\theoremstyle{definition}
\newtheorem{theorem}{Theorem}[section]
\newtheorem{lemma}[theorem]{Lemma}
\newtheorem{proposition}[theorem]{Proposition}
\newtheorem{corollary}[theorem]{Corollary}
\newtheorem{conjecture}[theorem]{Conjecture}
\newtheorem{definition}[theorem]{Definition}
\newtheorem{example}[theorem]{Example}
\newtheorem{notation}[theorem]{Notation}
\newtheorem{remark}[theorem]{Remark}

% Non-italic proof label
\makeatletter
\renewenvironment{proof}[1][\proofname]{\par
  \pushQED{\qed}%
  \normalfont \topsep6\p@\@plus6\p@\relax
  \trivlist
  \item[\hskip\labelsep\bfseries #1\@addpunct{.}]\ignorespaces
}{%
  \popQED\endtrivlist\@endpefalse
}
\makeatother

%%% WRAP ENVIRONMENTS IN COLORED BOXES %%%
%%% Theorem-like statements: NOT breakable - must appear whole on a single page %%%
%%% Proof / Remark / Notation: breakable %%%
\tcolorboxenvironment{definition}{
  enhanced,breakable=false,
  colback=defBG,colframe=defFR,
  boxrule=0.6pt,arc=2pt,
  left=3mm,right=3mm,top=1.5mm,bottom=1.5mm,
  before skip=5pt,after skip=5pt
}
\tcolorboxenvironment{theorem}{
  enhanced,breakable=false,
  colback=thmBG,colframe=thmFR,
  boxrule=0.6pt,arc=2pt,
  left=3mm,right=3mm,top=1.5mm,bottom=1.5mm,
  before skip=5pt,after skip=5pt
}
\tcolorboxenvironment{lemma}{
  enhanced,breakable=false,
  colback=lemBG,colframe=lemFR,
  boxrule=0.6pt,arc=2pt,
  left=3mm,right=3mm,top=1.5mm,bottom=1.5mm,
  before skip=5pt,after skip=5pt
}
\tcolorboxenvironment{corollary}{
  enhanced,breakable=false,
  colback=corBG,colframe=corFR,
  boxrule=0.6pt,arc=2pt,
  left=3mm,right=3mm,top=1.5mm,bottom=1.5mm,
  before skip=5pt,after skip=5pt
}
\tcolorboxenvironment{proposition}{
  enhanced,breakable=false,
  colback=propBG,colframe=propFR,
  boxrule=0.6pt,arc=2pt,
  left=3mm,right=3mm,top=1.5mm,bottom=1.5mm,
  before skip=5pt,after skip=5pt
}
\tcolorboxenvironment{remark}{
  enhanced,breakable,
  colback=remBG,colframe=remFR,
  leftrule=3pt,rightrule=0pt,toprule=0pt,bottomrule=0pt,
  arc=0pt,left=3mm,right=2mm,top=1mm,bottom=1mm,
  before skip=4pt,after skip=4pt
}
\tcolorboxenvironment{notation}{
  enhanced,breakable,
  colback=notBG,colframe=notFR,
  leftrule=3pt,rightrule=0pt,toprule=0pt,bottomrule=0pt,
  arc=0pt,left=3mm,right=2mm,top=1mm,bottom=1mm,
  before skip=4pt,after skip=4pt
}
\tcolorboxenvironment{proof}{
  enhanced,breakable,
  colback=proofBG,colframe=proofFR,
  leftrule=2pt,rightrule=0pt,toprule=0pt,bottomrule=0pt,
  arc=0pt,left=3mm,right=2mm,top=1mm,bottom=1mm,
  before skip=4pt,after skip=4pt
}

%%% CUSTOM COMMANDS (from gci.tex) %%%
\newcommand{\R}{\mathbb{R}}
\newcommand{\N}{\mathbb{N}}
\newcommand{\PP}{\mathsf{P}}
\newcommand{\EE}{\mathsf{E}}
\newcommand{\Var}{\operatorname{var}}
\newcommand{\Cov}{\operatorname{cov}}
\newcommand{\tr}{\operatorname{tr}}
\newcommand{\rank}{\operatorname{rank}}
\newcommand{\diag}{\operatorname{diag}}
\newcommand{\dd}{\mathrm{d}}
\newcommand{\ind}{\mathbf{1}}
\newcommand{\caln}{\mathcal{N}}
\newcommand{\Mn}{M_{n\times m}}

%%% SPACING %%%
\usepackage{parskip}
\setlength{\parskip}{3pt}
\setlength{\jot}{4pt}
\AtBeginEnvironment{definition}{\vspace{\parskip}}
\AtBeginEnvironment{lemma}{\vspace{\parskip}}
\AtBeginEnvironment{theorem}{\vspace{\parskip}}
\AtBeginEnvironment{corollary}{\vspace{\parskip}}
\AtBeginEnvironment{proposition}{\vspace{\parskip}}
\AtBeginEnvironment{notation}{\vspace{\parskip}}
\AtBeginEnvironment{remark}{\vspace{\parskip}}

%%% SECTION STYLING %%%
\usepackage{titlesec}
\titleformat{\section}
  {\color{hdrBG}\normalfont\Large\bfseries}{\thesection}{0.7em}{}
\titleformat{\subsection}
  {\color{accent}\normalfont\large\bfseries}{\thesubsection}{0.6em}{}
\titlespacing*{\section}{0pt}{4pt}{2pt}
\titlespacing*{\subsection}{0pt}{4pt}{1pt}

\title{\color{hdrBG}\bfseries\Huge The Gaussian Correlation Inequality after Royen}
\author{}
\date{}

\begin{document}

%%% TITLE PAGE %%%
\begin{titlepage}
\thispagestyle{empty}
\vspace*{\fill}
\begin{center}
{\color{hdrBG}\bfseries\fontsize{36}{44}\selectfont The Gaussian Correlation Inequality}\\[0.6em]
%{\color{hdrBG}\bfseries\fontsize{28}{34}\selectfont after Royen}\\[1.2em]
{\color{accent}\bfseries\Large Slide edition}\\[2em]
\rule{60mm}{0.8pt}\\[.5em]
{\large After T.\ Royen (2014) and R.\ Lata{\l}a \& D.\ Matlak (2015)}\\[.5em]
\rule{60mm}{0.8pt}
\end{center}
\vspace*{\fill}
\end{titlepage}


%% ====================================================================
\section{Introduction}\label{sec:intro}
%% ====================================================================

\begin{theorem}[Gaussian correlation inequality]\label{thm:gcc_geom}
  For any closed symmetric convex sets $K,L\subseteq\R^d$ and any
  centered Gaussian measure $\mu$ on $\R^d$,
  \begin{equation}\label{eq:gcc}
    \mu(K\cap L)\;\geqslant\;\mu(K)\,\mu(L).
  \end{equation}
\end{theorem}

These slides prove \Cref{thm:gcc_geom} by Royen's analytic argument,
following Lata{\l}a--Matlak's Gaussian-only exposition.  The inputs
are the Laplace transform of a squared Gaussian, a determinant
expansion, dominated convergence, and integration by parts.

\begin{notation}\label{not:conventions}
  \leavevmode
  \begin{itemize}[nosep]
    \item $\caln(0,C)$: centered Gaussian distribution with covariance
      $C$.
    \item $M_{n\times m}$: real $n\times m$ matrices.  $A\succ 0$
      (resp.\ $A\succcurlyeq 0$) means $A$ is symmetric positive
      (semi)definite.
    \item For $A=(a_{ij})_{i,j\leqslant n}$ and
      $J\subseteq[n]:=\{1,\ldots,n\}$,
      $A_J:=(a_{ij})_{i,j\in J}$ is the principal submatrix,
      $|A|$ the determinant, $|J|$ the cardinality.
    \item $I_n$: $n\times n$ identity.
    \item For $\varepsilon\in\{-1,1\}^n$, $x\in(0,\infty)^n$:
      $\sqrt{x}_\varepsilon:=(\varepsilon_i\sqrt{x_i})_{i\leqslant n}$.
  \end{itemize}
\end{notation}

We will use three standard facts without further comment.

\begin{lemma}[Square root]\label{lem:sqrt}
  For any $A\succcurlyeq 0$ there exists a unique $A^{1/2}\succcurlyeq 0$ with
  $(A^{1/2})^2=A$, and $A^{1/2}$ is invertible iff $A\succ 0$.
\end{lemma}

(Spectral theorem: diagonalize $A=O\Lambda O^\top$ with $\Lambda\succcurlyeq 0$,
set $A^{1/2}:=O\Lambda^{1/2}O^\top$; uniqueness is standard.)

\begin{lemma}[Sylvester determinant identity]\label{lem:sylvester}
  For $A_1\in M_{n\times m}$ and $A_2\in M_{m\times n}$,
  $|I_n+A_1A_2|=|I_m+A_2A_1|$.
\end{lemma}

\begin{proof}
  Compute $D:=\bigl|\begin{smallmatrix}I_n & -A_1\\ A_2 &
  I_m\end{smallmatrix}\bigr|$ two ways.  Row-reduce $R_1\leftarrow
  R_1+A_1R_2$ to clear the $(1,1)$-block's contribution to the $(1,2)$
  block: the matrix becomes
  $\bigl(\begin{smallmatrix}I_n+A_1A_2 & 0\\ A_2 & I_m\end{smallmatrix}\bigr)$,
  so $D=|I_n+A_1A_2|$.  Alternatively, column-reduce $C_1\leftarrow
  C_1+C_2A_2$: the matrix becomes
  $\bigl(\begin{smallmatrix}I_n & -A_1\\ 0 & I_m+A_2A_1\end{smallmatrix}\bigr)$,
  so $D=|I_m+A_2A_1|$.
\end{proof}

\begin{lemma}[Strict separation; Hahn--Banach]\label{lem:separation}
  Let $K_1,K_2\subseteq\R^d$ be disjoint closed convex sets with $K_1$
  bounded.  There exist $w\in\R^d\setminus\{0\}$ and $\alpha_1<\alpha_2$
  in $\R$ with $\langle w,x\rangle\leqslant\alpha_1$ for $x\in K_1$ and
  $\langle w,x\rangle\geqslant\alpha_2$ for $x\in K_2$.
\end{lemma}

(Standard corollary of Hahn--Banach in finite dimension; see e.g.\
Rockafellar, \emph{Convex Analysis}, Cor.\ 11.4.2.)


%% ====================================================================
\section{Setup and reduction to rectangles}\label{sec:setup}
%% ====================================================================

\begin{definition}[Centered Gaussian measure]\label{def:gaussian}
  For $C\succ 0$, the measure $\caln(0,C)$ on $\R^n$ has density
  \[
    \phi_C(x)
    \;=\;
    (2\pi)^{-n/2}|C|^{-1/2}
    \exp\bigl(-\tfrac{1}{2}\langle C^{-1}x,x\rangle\bigr).
  \]
\end{definition}

We assume $C\succ 0$ throughout the proof.

\begin{definition}\label{def:symconvex}
  $K\subseteq\R^n$ is \emph{symmetric} if $K=-K$.  A \emph{symmetric
  slab} is
  $S_{v,t}:=\{x\in\R^n:|\langle v,x\rangle|\leqslant t\}$ for
  $v\in\R^n\setminus\{0\}$, $t>0$.
\end{definition}

\begin{lemma}[Slab approximation]\label{lem:slab_approx}
  Every closed symmetric convex set $K\subseteq\R^d$ is a countable
  intersection of symmetric slabs $S_{v,t}$.
\end{lemma}

\begin{proof}
  \leavevmode
  \begin{enumerate}
    \item Reduction to bounded $K$.  Set $K_m:=K\cap[-m,m]^d$, a closed
      symmetric convex set.  Since $K=\bigcup_m K_m$, by monotone
      convergence $\mu(K)=\lim_m\mu(K_m)$ and similarly for $K\cap L$,
      so it suffices to prove the slab approximation for bounded $K$;
      we assume $K\subseteq B(0,R)$ for some $R<\infty$ below.

    \item Bounded case.  Let
      $\mathcal V:=\{(v,t)\in\mathbb{Q}^d\times\mathbb{Q}_{>0}:K\subseteq S_{v,t}\}$,
      countable.  Trivially $K\subseteq\bigcap_{(v,t)\in\mathcal V}S_{v,t}$.

      Conversely, fix $y\notin K$ and let $\varepsilon>0$ with
      $B(y,\varepsilon)\cap K=\varnothing$.  Applying
      \Cref{lem:separation} to $K_1:=\overline{B(y,\varepsilon/2)}$
      (bounded) and $K_2:=K$ (closed convex, disjoint from $K_1$) yields
      $w\in\R^d\setminus\{0\}$ and $a\in\R$ with
      $\sup_{x\in K}\langle w,x\rangle\leqslant a$ and
      $\langle w,y\rangle\geqslant a+\varepsilon|w|$.
      Since $K$ is symmetric and convex, $0=\tfrac12 x+\tfrac12(-x)\in K$
      for any $x\in K$, hence $a\geqslant\langle w,0\rangle=0$; and
      $K=-K$ gives $\sup_{x\in K}|\langle w,x\rangle|\leqslant a$.
      Pick $(w',a')\in\mathbb{Q}^d\times\mathbb{Q}_{>0}$ with
      $|w'-w|<\delta$, $|a'-a|<\delta$.  Then
      \[
        \sup_{x\in K}|\langle w',x\rangle|
        \leqslant a+\delta R,\qquad
        |\langle w',y\rangle|\geqslant a+\varepsilon|w|-\delta|y|.
      \]
      Choosing $\delta<\min(\varepsilon|w|/(|y|+R+1),\varepsilon|w|/2)$
      and $a':=a+\delta R+\delta/2\in\mathbb{Q}_{>0}$ gives
      $K\subseteq S_{w',a'}$ and $y\notin S_{w',a'}$.
  \end{enumerate}
\end{proof}

By \Cref{lem:slab_approx} and monotone convergence, the GCI reduces
to the case where $K$ and $L$ are finite intersections of symmetric
slabs.  Writing
\[
  K=\bigl\{x : |\langle v_i,x\rangle|\leqslant t_i,\;1\leqslant i\leqslant n_1\bigr\},
  \qquad
  L=\bigl\{x : |\langle v_i,x\rangle|\leqslant t_i,\;n_1{+}1\leqslant i\leqslant n\bigr\},
\]
setting $n=n_1+n_2$ and $X_i:=\langle v_i,G\rangle$ for $G\sim\mu$,
the GCI reduces to:

\begin{theorem}[GCI, analytic form]\label{thm:gcc2}
  Let $X=(X_1,\ldots,X_n)$ be centered Gaussian with $C\succ 0$,
  $n=n_1+n_2$.  For any $t_1,\ldots,t_n>0$,
  \begin{equation}\label{eq:gcc2}
    \PP(|X_i|\leqslant t_i,\;i\leqslant n)
    \;\geqslant\;
    \PP(|X_i|\leqslant t_i,\;i\leqslant n_1)\,
    \PP(|X_i|\leqslant t_i,\;n_1{<}i\leqslant n).
  \end{equation}
\end{theorem}


%% ====================================================================
\section{Auxiliary lemmas}\label{sec:lemmas}
%% ====================================================================

\subsection{Laplace transform of squared Gaussian vectors}

\begin{lemma}\label{lem:Lap}
  For $X\sim\caln(0,C)$ in $\R^n$ and $\lambda_i\geqslant 0$,
  \[
    \EE\exp\Bigl(-\sum_{i=1}^n\lambda_iX_i^2\Bigr)
    \;=\;|I_n+2\Lambda C|^{-1/2},
    \qquad\Lambda:=\diag(\lambda_1,\ldots,\lambda_n).
  \]
\end{lemma}

\begin{proof}
  \leavevmode
  \begin{enumerate}
    \item Standard normal quadratic.  For $Y\sim\caln(0,I_n)$ and a
      symmetric matrix $B$ with $I_n-2B\succ 0$, combine the density of
      $Y$ with the factor $e^{\langle BY,Y\rangle}$: the exponent is
      $-\tfrac12\langle y,y\rangle+\langle By,y\rangle
      =-\tfrac12\langle(I_n-2B)y,y\rangle$, so
      \begin{equation}\label{eq:gauss_quad}
        \EE e^{\langle BY,Y\rangle}
        =(2\pi)^{-n/2}\int_{\R^n}e^{-\frac12\langle(I_n-2B)y,y\rangle}\,\dd y
        =|I_n-2B|^{-1/2},
      \end{equation}
      using $\int_{\R^n}e^{-\frac12\langle Ay,y\rangle}\dd y=(2\pi)^{n/2}|A|^{-1/2}$
      (change variables $y=A^{-1/2}z$; $A=I_n-2B$).

    \item Reduction to the standard case.  Set $A:=C^{1/2}$
      (\Cref{lem:sqrt}); then $C=AA^\top$ and $X:=AY\sim\caln(0,C)$.
      Taking $B:=-A^\top\Lambda A$ in \eqref{eq:gauss_quad}: since
      $\lambda_i\geqslant 0$ gives $\Lambda\succcurlyeq 0$, hence $B\preccurlyeq 0$
      and $I_n-2B\succcurlyeq I_n\succ 0$.  The exponent becomes
      $\langle BY,Y\rangle=-\langle\Lambda AY,AY\rangle
      =-\sum_i\lambda_i(AY)_i^2=-\sum_i\lambda_iX_i^2$, so
      \[
        \EE\exp\Bigl(-\sum_i\lambda_iX_i^2\Bigr)
        =|I_n+2A^\top\Lambda A|^{-1/2}.
      \]

    \item Commuting the factors.  Apply \Cref{lem:sylvester} with
      $A_1:=A^\top$, $A_2:=2\Lambda A$:
      $|I_n+2A^\top\Lambda A|=|I_n+2\Lambda AA^\top|=|I_n+2\Lambda C|$.
  \end{enumerate}
\end{proof}


\subsection{A determinant identity}

\begin{lemma}\label{lem:det}
  \leavevmode
  \begin{enumerate}[label=(\roman*),nosep]
    \item For $A\in M_{n\times n}$,
    \begin{equation}\label{eq:det1}
      |I_n+A|\;=\;1+\sum_{\varnothing\ne J\subseteq[n]}|A_J|.
    \end{equation}
    \item For $A=\bigl(\begin{smallmatrix}A_{11}&A_{12}\\A_{21}&A_{22}
      \end{smallmatrix}\bigr)$ with $A_{ii}$ invertible,
    \begin{equation}\label{eq:det2}
      |A|=|A_{11}||A_{22}|\cdot
      \bigl|I_{n_1}-A_{11}^{-1/2}A_{12}A_{22}^{-1}A_{21}A_{11}^{-1/2}\bigr|.
    \end{equation}
    If $A\succ 0$, then $0\leqslant A_{11}^{-1/2}A_{12}A_{22}^{-1}A_{21}
    A_{11}^{-1/2}\leqslant I_{n_1}$.
  \end{enumerate}
\end{lemma}

\begin{proof}
  \leavevmode
  \begin{enumerate}[label=(\roman*)]
    \item By Leibniz,
      $|I_n+A|=\sum_\sigma\mathrm{sgn}(\sigma)\prod_i
      (\delta_{i\sigma(i)}+a_{i\sigma(i)})$.  Expanding the product, each
      summand corresponds to a choice of $J\subseteq[n]$ (positions where
      $a_{i\sigma(i)}$ is selected); non-zero contributions require
      $\sigma(i)=i$ for $i\notin J$.  The $\sigma$'s fixing $[n]\setminus J$
      correspond to permutations of $J$, giving $\sum_J|A_J|$; the
      $J=\varnothing$ term is $1$.

    \item Let $A_{11}^{1/2},A_{22}^{1/2}$ be the positive symmetric
      square roots (\Cref{lem:sqrt}), invertible by assumption.
      Setting $P:=A_{11}^{-1/2}A_{12}A_{22}^{-1/2}$ and
      $Q:=A_{22}^{-1/2}A_{21}A_{11}^{-1/2}$, a direct block
      multiplication gives
      \[
        A=\diag(A_{11}^{1/2},A_{22}^{1/2})\,
        \bigl(\begin{smallmatrix}I&P\\Q&I\end{smallmatrix}\bigr)\,
        \diag(A_{11}^{1/2},A_{22}^{1/2}),
      \]
      so $|A|=|A_{11}||A_{22}|\cdot\bigl|
      \begin{smallmatrix}I&P\\Q&I\end{smallmatrix}\bigr|$.
      Row-reduce by $R_1\leftarrow R_1-PR_2$ (determinant invariant):
      \[
        \bigl(\begin{smallmatrix}I&P\\Q&I\end{smallmatrix}\bigr)
        \;\to\;
        \bigl(\begin{smallmatrix}I-PQ&0\\Q&I\end{smallmatrix}\bigr),
      \]
      block-triangular with determinant $|I-PQ|$.  Substituting
      $PQ=A_{11}^{-1/2}A_{12}A_{22}^{-1}A_{21}A_{11}^{-1/2}$ gives
      \eqref{eq:det2}.

      When $A\succ 0$: for all $s,t\in\R$ and $x\in\R^{n_1}$,
      $y\in\R^{n_2}$,
      \[
        \langle A(sx,ty)^\top,(sx,ty)^\top\rangle
        =s^2\langle A_{11}x,x\rangle+2st\langle A_{21}x,y\rangle
        +t^2\langle A_{22}y,y\rangle\geqslant 0.
      \]
      Nonnegativity of this quadratic in $(s,t)$ forces
      $\langle A_{21}x,y\rangle^2\leqslant\langle A_{11}x,x\rangle
      \langle A_{22}y,y\rangle$.
      Setting $x:=A_{11}^{-1/2}u$, $y:=A_{22}^{-1/2}v$:
      $\langle Qu,v\rangle^2\leqslant|u|^2|v|^2$, i.e.,
      $\|Q\|_{\mathrm{op}}\leqslant 1$, equivalently
      $0\preccurlyeq Q^\top Q\preccurlyeq I_{n_1}$.  For $A$ symmetric,
      $Q=P^\top$, so
      $A_{11}^{-1/2}A_{12}A_{22}^{-1}A_{21}A_{11}^{-1/2}=PP^\top=Q^\top Q$.
  \end{enumerate}
\end{proof}


\subsection{Density of $Z(\tau)$ and the interchange lemma}

Partition
\begin{equation}\label{eq:partition}
  C=\begin{pmatrix}C_{11}&C_{12}\\C_{21}&C_{22}\end{pmatrix},
  \quad C_{ij}\in M_{n_i\times n_j},\;n_1+n_2=n,
\end{equation}
and define the interpolating family
\begin{equation}\label{eq:Ctau}
  C(\tau):=\begin{pmatrix}C_{11}&\tau C_{12}\\
    \tau C_{21}&C_{22}\end{pmatrix}
  =\tau C(1)+(1-\tau)C(0),\qquad\tau\in[0,1].
\end{equation}
Since $C(0)=\diag(C_{11},C_{22})\succ 0$ and $C(1)=C\succ 0$, each
$C(\tau)\succ 0$ by convexity.  Let $X(\tau)\sim\caln(0,C(\tau))$ and
$Z_i(\tau):=\tfrac{1}{2}X_i(\tau)^2$.

\begin{lemma}[Density and interchange]\label{lem:diff}
  Let $f(x,\tau)$ denote the density of $Z(\tau)$ on $(0,\infty)^n$.
  \begin{enumerate}[label=(\roman*),nosep]
    \item
    \begin{equation}\label{eq:fxtau}
      f(x,\tau)
      =|C(\tau)|^{-1/2}(4\pi)^{-n/2}
      \frac{1}{\sqrt{x_1\cdots x_n}}
      \sum_{\varepsilon\in\{-1,1\}^n}
      e^{-\langle C(\tau)^{-1}\sqrt{x}_\varepsilon,
        \sqrt{x}_\varepsilon\rangle}
      \ind_{(0,\infty)^n}(x).
    \end{equation}
    \item For any Borel $K\subseteq[0,\infty)^n$ and
    $\lambda_i\geqslant 0$,
    \begin{equation}\label{eq:interchange}
      \int_K e^{-\sum_i\lambda_ix_i}\,
      \frac{\partial}{\partial\tau}f(x,\tau)\,\dd x
      \;=\;
      \frac{\partial}{\partial\tau}\int_K e^{-\sum_i\lambda_ix_i}\,
      f(x,\tau)\,\dd x.
    \end{equation}
  \end{enumerate}
\end{lemma}

\begin{proof}
  \leavevmode
  \begin{enumerate}[label=(\roman*)]
    \item The density of $X(\tau)$ is $\phi_{C(\tau)}$ from
      \Cref{def:gaussian}.  On each sign orthant
      $\R^n_\varepsilon:=\{y:\text{sgn}(y_i)=\varepsilon_i\}$, the map
      $y\mapsto x_i:=\tfrac12 y_i^2$ is a diffeomorphism onto
      $(0,\infty)^n$ with inverse $y_i=\varepsilon_i\sqrt{2x_i}$ and
      Jacobian $\dd y_i/\dd x_i=\varepsilon_i/\sqrt{2x_i}$; absolute
      Jacobian $\prod_i(2x_i)^{-1/2}$.  Summing over the $2^n$ sign
      patterns and using
      $\langle C(\tau)^{-1}y,y\rangle
      =2\langle C(\tau)^{-1}\sqrt{x}_\varepsilon,\sqrt{x}_\varepsilon\rangle$
      (by $y_i=\varepsilon_i\sqrt{2x_i}$):
      \begin{align*}
        f(x,\tau)
        &=\sum_\varepsilon\phi_{C(\tau)}(\sqrt{2x}_\varepsilon)
          \cdot\prod_i(2x_i)^{-1/2}\\
        &=|C(\tau)|^{-1/2}(2\pi)^{-n/2}2^{-n/2}(x_1\cdots x_n)^{-1/2}
          \sum_\varepsilon e^{-\langle C(\tau)^{-1}\sqrt{x}_\varepsilon,
          \sqrt{x}_\varepsilon\rangle},
      \end{align*}
      which is \eqref{eq:fxtau} after combining
      $(2\pi)^{-n/2}2^{-n/2}=(4\pi)^{-n/2}$.

    \item We build a $\tau$-uniform integrable dominating function for
      $\partial_\tau f$ via four steps.
      \begin{enumerate}
        \item Eigenvalue lower bound.  Each $C(\tau)\succ 0$, and
          $\tau\mapsto\lambda_{\min}(C(\tau)^{-1})=1/\lambda_{\max}(C(\tau))$
          is continuous and positive on compact $[0,1]$, so bounded
          below by some $a>0$.  Then
          $\langle C(\tau)^{-1}\sqrt{x}_\varepsilon,\sqrt{x}_\varepsilon\rangle
          \geqslant a|\sqrt{x}_\varepsilon|^2=a\sum_ix_i$.

        \item Derivative formula.  Differentiating
          $C(\tau)C(\tau)^{-1}=I_n$,
          $\tfrac{\dd}{\dd\tau}C(\tau)^{-1}=-C(\tau)^{-1}C'(\tau)C(\tau)^{-1}$,
          with $C'(\tau)=C(1)-C(0)$ constant.  Since
          $\tau\mapsto C(\tau)^{-1}$ is continuous on compact $[0,1]$,
          both $|C(\tau)|^{-1/2}$ and
          $\|\tfrac{\dd}{\dd\tau}C(\tau)^{-1}\|_{\mathrm{op}}$ are
          bounded, say by $M,b<\infty$.  Hence
          $|\partial_\tau|C(\tau)|^{-1/2}|\leqslant M$ and
          $|\langle\tfrac{\dd}{\dd\tau}C(\tau)^{-1}\sqrt{x}_\varepsilon,
          \sqrt{x}_\varepsilon\rangle|\leqslant b|\sqrt{x}_\varepsilon|^2
          =b\sum_ix_i$.

        \item Dominating function.  Differentiating \eqref{eq:fxtau}
          in $\tau$ (product rule on the prefactor $|C(\tau)|^{-1/2}$
          and the exponent), each of the $2^n$ sign-orthant summands
          is bounded by
          \[
            \bigl(M+b\textstyle\sum_ix_i\bigr)
            (4\pi)^{-n/2}(x_1\cdots x_n)^{-1/2}e^{-a\sum_ix_i}.
          \]
          Summing over $\varepsilon$ and absorbing constants,
          \[
            \sup_{\tau\in[0,1]}|\partial_\tau f(x,\tau)|
            \leqslant g(x):=c\,\frac{1+\sum_ix_i}{\sqrt{x_1\cdots x_n}}
            e^{-a\sum_ix_i}
          \]
          for some $c=c(C,n)$.  Fubini:
          $\int_0^\infty x^{-1/2}(1+x)e^{-ax}\,\dd x<\infty$, so
          factoring in each coordinate gives $\int g<\infty$.

        \item Interchange.  For any Borel $K\subseteq[0,\infty)^n$ and
          $\lambda_i\geqslant 0$,
          $e^{-\sum_i\lambda_ix_i}|\partial_\tau f|\leqslant g$
          uniformly in $\tau$.  Dominated convergence applied to the
          difference quotient
          $(f(x,\tau+h)-f(x,\tau))/h\to\partial_\tau f(x,\tau)$ gives
          \eqref{eq:interchange}.
      \end{enumerate}
  \end{enumerate}
\end{proof}


\subsection{The noncentral gamma and the $h$-family}

\begin{definition}\label{def:ncgamma}
  For $\alpha>0$, $y\geqslant 0$, the \emph{noncentral gamma density} is
  \[
    g_\alpha(x,y)
    :=e^{-x-y}\sum_{k=0}^\infty
    \frac{x^{k+\alpha-1}}{\Gamma(k+\alpha)}\,\frac{y^k}{k!},
    \qquad x>0.
  \]
\end{definition}

\begin{remark}\label{rem:ncgamma_poisson}
  Equivalently,
  \[
    g_\alpha(x,y)=\sum_{k\geqslant 0}e^{-y}\tfrac{y^k}{k!}g_{\alpha+k}(x),
  \]
  a Poisson$(y)$ mixture of central gamma$(\alpha+k)$ densities; so
  $g_\alpha(\cdot,y)$ is a probability density for each $y\geqslant 0$.
\end{remark}

\begin{definition}\label{def:h_family}
  For $\mu>0$, $\alpha\in(0,\infty)^n$, and
  $Y=(Y_1,\ldots,Y_n)$ with $Y_i\geqslant 0$ a.s.,
  \[
    h_{\alpha,\mu,Y}(x)
    :=\EE\prod_{i=1}^n
      \tfrac{1}{\mu}\,g_{\alpha_i}\!\bigl(\tfrac{x_i}{\mu},Y_i\bigr),
    \qquad x\in(0,\infty)^n.
  \]
\end{definition}

\begin{lemma}[Properties of the $h$-family]\label{lem:proph}
  Write $h_\alpha:=h_{\alpha,\mu,Y}$.
  \begin{enumerate}[label=(\roman*)]
    \item $h_\alpha\geqslant 0$ and $\int h_\alpha\,\dd x=1$.
    \item If $\alpha_i>1$, then $h_\alpha(x)\to 0$ as $x_i\to 0+$,
      and
      \begin{equation}\label{eq:deriv_h}
        \frac{\partial}{\partial x_i}h_\alpha=h_{\alpha-e_i}-h_\alpha.
      \end{equation}
    \item If $\alpha\in(1,\infty)^n$, then for any $J\subseteq[n]$
      and $\lambda_i\geqslant 0$,
      $\partial^{|J|}h_\alpha/\partial x_J\in L_1((0,\infty)^n)$ and
      \begin{equation}\label{eq:proph_iii}
        \int e^{-\sum_i\lambda_ix_i}
          \frac{\partial^{|J|}}{\partial x_J}h_\alpha\,\dd x
        \;=\;
        \Bigl(\prod_{i\in J}\lambda_i\Bigr)
        \int e^{-\sum_i\lambda_ix_i}h_\alpha\,\dd x.
      \end{equation}
  \end{enumerate}
\end{lemma}

\begin{proof}
  \leavevmode
  \begin{enumerate}[label=(\roman*)]
    \item $\int_0^\infty\mu^{-1}g_\alpha(x_i/\mu,Y_i)\,\dd x_i
      =\int_0^\infty g_\alpha(u,Y_i)\,\dd u=1$ for each $\omega$
      (\Cref{rem:ncgamma_poisson}); Fubini gives
      $\int h_\alpha=\EE[1]=1$.

    \item We split into two parts.
      \begin{enumerate}
        \item Gamma bound.  $\log\Gamma$ is strictly convex on
          $(0,\infty)$, and since $\Gamma(1)=\Gamma(2)=1$, $\Gamma$ has
          a unique minimum at some $s_0\in(1,2)$ with $\Gamma(s_0)>\tfrac12$
          (numerically $s_0\approx 1.462$, $\Gamma(s_0)\approx 0.886$).
          Hence
          \begin{equation}\label{eq:gamma_lb}
            \Gamma(s)\geqslant\tfrac12\quad\text{for all }s>0,
          \end{equation}
          giving $\Gamma(\alpha)^{-1}\leqslant 2$ for the $k=0$ term.
          For $k\geqslant 1$, induction via
          $\Gamma(k+\alpha)=(k-1+\alpha)\Gamma(k-1+\alpha)\geqslant
          k\cdot\tfrac12(k-2)!=\tfrac12(k-1)!$ gives
          $\Gamma(k+\alpha)\geqslant\tfrac12(k-1)!$ (base case $k=1$:
          $\Gamma(1+\alpha)\geqslant\tfrac12$).  Substituting into
          $g_\alpha(x,y)$:
          \begin{align*}
            g_\alpha(x,y)
            &\leqslant e^{-x-y}\Bigl[2x^{\alpha-1}
              +\sum_{k\geqslant 1}\tfrac{2x^{k+\alpha-1}y^k}{(k-1)!k!}\Bigr]\\
            &=2x^{\alpha-1}e^{-x-y}
              +2x^\alpha e^{-x}\underbrace{\sum_{k\geqslant 1}\tfrac{x^{k-1}}{(k-1)!}
              \cdot\tfrac{y^ke^{-y}}{k!}}_{\leqslant\;e^x\cdot 1}
            \leqslant 2x^{\alpha-1}(e^{-x}+x),
          \end{align*}
          using $y^ke^{-y}/k!\leqslant\sup_y(y^ke^{-y}/k!)\leqslant 1$
          (term in a Poisson distribution), and
          $\sum_{k\geqslant 1}x^{k-1}/(k-1)!=e^x$.  Thus
          \begin{equation}\label{eq:gbound}
            g_\alpha(x,y)\leqslant 2x^{\alpha-1}(e^{-x}+x),
          \end{equation}
          and substituting into \Cref{def:h_family} (note $Y_i$ drops out),
          \begin{equation}\label{eq:hbound}
            h_\alpha(x)\leqslant(2/\mu)^n\prod_{i}(x_i/\mu)^{\alpha_i-1}(1+x_i/\mu).
          \end{equation}
          When $\alpha_i>1$, the factor $(x_i/\mu)^{\alpha_i-1}\to 0$
          as $x_i\to 0+$, so $h_\alpha(x)\to 0$.

        \item Derivative formula.  Differentiating the series in
          \Cref{def:ncgamma} term by term (using
          $\Gamma(k+\alpha)=(k+\alpha-1)\Gamma(k+\alpha-1)$, valid when
          $\alpha>0$),
          \begin{equation}\label{eq:g_deriv}
            \partial_x g_\alpha(x,y)=g_{\alpha-1}(x,y)-g_\alpha(x,y),
            \qquad\alpha>1,
          \end{equation}
          with convergence locally uniform by \eqref{eq:gbound} applied
          to $g_{\alpha-1}$ (valid since $\alpha-1>0$).  Inserting into
          \Cref{def:h_family}: the integrand and its $x_i$-derivative are
          bounded by an integrable function of $(x,\omega)$ uniformly in
          $x_i$ on compact subsets of $(0,\infty)$ (use \eqref{eq:gbound}
          for $g_{\alpha_i-1}(\cdot,Y_i)$ and $g_{\alpha_i}(\cdot,Y_i)$),
          justifying dominated convergence to differentiate under $\EE$,
          which gives \eqref{eq:deriv_h}.
      \end{enumerate}

    \item Iterating \eqref{eq:deriv_h},
      \begin{equation}\label{eq:derivJ_h}
        \frac{\partial^{|J|}}{\partial x_J}h_\alpha
        =\sum_{\delta\in\{0,1\}^J}(-1)^{|J|-|\delta|}
          h_{\alpha-\sum_{i\in J}\delta_ie_i},
      \end{equation}
      each summand in $L_1$ by (i) (all $\alpha_i-\delta_i\geqslant\alpha_i-1>0$).
      Integration by parts in each $x_j$, $j\in J$: the boundary term at
      $x_j\to\infty$ vanishes by \eqref{eq:hbound} (exponential decay in
      every coordinate via $\int_0^\infty x^{\alpha-1}(1+x)e^{-ax}\,\dd x<\infty$,
      implying pointwise decay of the partial integral); at $x_j\to 0+$
      by (ii), applied to the $\alpha'_j:=\alpha_j-\delta_j\geqslant\alpha_j-1>0$
      shape, together with the extra factor $e^{-\sum_i\lambda_ix_i}$
      which is bounded.  Each integration by parts brings down a
      $\lambda_j$, giving \eqref{eq:proph_iii}.
  \end{enumerate}
\end{proof}


\subsection{The family $h_{k,C}$ and its Laplace transform}
  \label{sec:hkC}

\begin{definition}\label{def:hkC}
  For $C\succ 0$ and $\mu>0$ with $C-\mu I_n\succ 0$, set
  $A:=(C-\mu I_n)^{1/2}$ (\Cref{lem:sqrt}, which is well-defined and
  invertible), so that $C=\mu I_n+AA^\top$.  Let $(g_j^{(l)})$,
  $j\leqslant n$, $l\leqslant k$, be i.i.d.\ $\caln(0,1)$, and
  \begin{equation}\label{eq:defhkC}
    Y_i:=\sum_{l=1}^k\Bigl(\sum_{j=1}^n
      \tfrac{1}{\sqrt{2\mu}}g_j^{(l)}a_{ij}\Bigr)^2,
    \qquad
    h_{k,C}:=h_{(k/2,\ldots,k/2),\mu,Y}.
  \end{equation}
\end{definition}

\begin{lemma}\label{lem:LaplhkC}
  For $\lambda_i\geqslant 0$,
  $\int e^{-\sum_i\lambda_ix_i}h_{k,C}\,\dd x
  =|I_n+\Lambda C|^{-k/2}$.
\end{lemma}

\begin{proof}
  \leavevmode
  \begin{enumerate}
    \item Scalar Laplace transform.  Substitute $u:=x/\mu$ and expand
      the series in \Cref{def:ncgamma}:
      \begin{align*}
        \int_0^\infty\mu^{-1}e^{-\lambda x}g_\alpha(x/\mu,y)\,\dd x
        &=\int_0^\infty e^{-\lambda\mu u-u-y}
          \sum_{k\geqslant 0}\tfrac{u^{k+\alpha-1}}{\Gamma(k+\alpha)}
          \tfrac{y^k}{k!}\,\dd u\\
        &=e^{-y}\sum_{k\geqslant 0}\tfrac{y^k}{k!}(1+\mu\lambda)^{-(k+\alpha)}\\
        &=(1+\mu\lambda)^{-\alpha}\exp\bigl(-\tfrac{\mu\lambda}{1+\mu\lambda}y\bigr),
      \end{align*}
      using $\int_0^\infty u^{k+\alpha-1}e^{-(1+\mu\lambda)u}\,\dd u
      =\Gamma(k+\alpha)(1+\mu\lambda)^{-(k+\alpha)}$ (Fubini-term-by-term
      justified by absolute convergence).  So
      \begin{equation}\label{eq:1d_lt}
        \int_0^\infty\mu^{-1}e^{-\lambda x}g_\alpha(x/\mu,y)\,\dd x
        =(1+\mu\lambda)^{-\alpha}\exp\bigl(-\tfrac{\mu\lambda}{1+\mu\lambda}y\bigr).
      \end{equation}

    \item Product form.  Applying \eqref{eq:1d_lt} coordinatewise with
      $\alpha_i=k/2$ and Fubini (nonneg integrand),
      \begin{equation}\label{eq:hkC_step1}
        \int e^{-\sum_i\lambda_ix_i}h_{k,C}\,\dd x
        =\prod_i(1+\mu\lambda_i)^{-k/2}\cdot
        \EE\exp\Bigl(-\sum_i\tfrac{\mu\lambda_i}{1+\mu\lambda_i}Y_i\Bigr).
      \end{equation}

    \item Evaluate the expectation.  From \eqref{eq:defhkC}, set
      $X^{(l)}:=\tfrac{1}{\sqrt{2\mu}}Ag^{(l)}$, so
      $X^{(l)}\sim\caln(0,\tfrac{1}{2\mu}AA^\top)$ i.i.d.\ across $l$
      and $Y_i=\sum_{l=1}^k(X_i^{(l)})^2$.  By independence across $l$
      and \Cref{lem:Lap} applied to each $X^{(l)}$ with
      $\tilde\Lambda:=\diag(\tfrac{\mu\lambda_i}{1+\mu\lambda_i})$,
      \[
        \EE\exp\Bigl(-\sum_i\tfrac{\mu\lambda_i}{1+\mu\lambda_i}Y_i\Bigr)
        =\Bigl(\EE\exp(-\langle\tilde\Lambda X^{(1)},X^{(1)}\rangle)\Bigr)^k
        =|I_n+\tfrac{1}{\mu}\tilde\Lambda AA^\top|^{-k/2}.
      \]

    \item Combine.  Inserting step~3 into \eqref{eq:hkC_step1} and
      using $\prod_i(1+\mu\lambda_i)^{-k/2}=|I_n+\mu\Lambda|^{-k/2}$
      (since $I_n+\mu\Lambda$ is diagonal),
      \[
        \int e^{-\sum_i\lambda_ix_i}h_{k,C}\,\dd x
        =\bigl(|I_n+\mu\Lambda|\cdot
          |I_n+\tfrac{1}{\mu}\tilde\Lambda AA^\top|\bigr)^{-k/2}.
      \]
      The diagonal identity
      $\tilde\Lambda_i=\mu\lambda_i/(1+\mu\lambda_i)$ rearranges to
      $(I_n+\mu\Lambda)\tilde\Lambda=\mu\Lambda$.  Hence, using the
      multiplicativity of $\det$,
      \begin{align*}
        |I_n+\mu\Lambda|\cdot|I_n+\tfrac{1}{\mu}\tilde\Lambda AA^\top|
        &=|(I_n+\mu\Lambda)(I_n+\tfrac{1}{\mu}\tilde\Lambda AA^\top)|\\
        &=|I_n+\mu\Lambda
          +\tfrac{1}{\mu}(I_n+\mu\Lambda)\tilde\Lambda\,AA^\top|\\
        &=|I_n+\mu\Lambda+\Lambda AA^\top|
        =|I_n+\Lambda(\mu I_n+AA^\top)|=|I_n+\Lambda C|,
      \end{align*}
      giving the claim.
  \end{enumerate}
\end{proof}


\subsection{Uniqueness of the Laplace transform}

\begin{theorem}\label{thm:laplace_inj}
  If $\varphi_1,\varphi_2\in L_1((0,\infty)^n)$ satisfy
  $\int e^{-\sum_i\lambda_ix_i}\varphi_1\,\dd x
  =\int e^{-\sum_i\lambda_ix_i}\varphi_2\,\dd x$ for all
  $\lambda\in[0,\infty)^n$, then $\varphi_1=\varphi_2$ a.e.
\end{theorem}

\begin{proof}
  Set $\varphi:=\varphi_1-\varphi_2\in L_1((0,\infty)^n)$; the
  hypothesis becomes $\int e^{-\sum_i\lambda_ix_i}\varphi\,\dd x=0$
  for all $\lambda\in[0,\infty)^n$.
  \begin{enumerate}
    \item Change of variables.  With $y_i:=e^{-x_i}$, the map
      $(0,\infty)^n\to(0,1]^n$ has Jacobian $\prod_iy_i^{-1}$, so
      setting
      $\tilde\varphi(y):=\varphi(-\log y)\prod_iy_i^{-1}\in L_1((0,1]^n)$
      and noting $e^{-\lambda_ix_i}=y_i^{\lambda_i}$,
      \begin{equation}\label{eq:moment_zero}
        \int_{(0,1]^n}y^\lambda\tilde\varphi(y)\,\dd y=0
        \qquad\text{for all }\lambda\in[0,\infty)^n.
      \end{equation}

    \item From polynomials to continuous functions.  Restricting
      \eqref{eq:moment_zero} to $\lambda\in\mathbb{Z}_{\geqslant 0}^n$
      and taking finite linear combinations,
      $\int_{[0,1]^n}P(y)\tilde\varphi(y)\,\dd y=0$ for every polynomial
      $P$ (extend $\tilde\varphi$ by zero at $y_i=0$).  By
      Stone--Weierstrass, polynomials are uniformly dense in
      $C([0,1]^n)$; given $\psi\in C([0,1]^n)$ and $\varepsilon>0$,
      pick $P$ with $\|\psi-P\|_\infty<\varepsilon$, so
      \[
        \Bigl|\int\psi\tilde\varphi\Bigr|
        =\Bigl|\int(\psi-P)\tilde\varphi\Bigr|
        \leqslant\varepsilon\|\tilde\varphi\|_1.
      \]
      Sending $\varepsilon\to 0$:
      $\int_{[0,1]^n}\psi\tilde\varphi\,\dd y=0$ for every
      $\psi\in C([0,1]^n)$.

    \item From continuous functions to $\tilde\varphi=0$ a.e.  Let
      $E:=\{\tilde\varphi>0\}$.  By inner/outer regularity of Lebesgue
      measure, there exist compact $K_m\subseteq E$ and open
      $U_m\supseteq E$ with $|U_m\setminus K_m|\to 0$; Urysohn gives
      $\psi_m\in C([0,1]^n)$ with
      $\mathbf{1}_{K_m}\leqslant\psi_m\leqslant\mathbf{1}_{U_m}$.  Then
      $\psi_m\to\mathbf{1}_E$ a.e.\ along a subsequence, and
      $|\psi_m\tilde\varphi|\leqslant|\tilde\varphi|\in L_1$, so
      dominated convergence gives
      $\int_E\tilde\varphi=\lim\int\psi_m\tilde\varphi=0$.  Similarly
      for $\{\tilde\varphi<0\}$, so $\int|\tilde\varphi|=0$ and
      $\tilde\varphi=0$ a.e., hence $\varphi=0$ a.e.
  \end{enumerate}
\end{proof}


%% ====================================================================
\section{Proof of \Cref{thm:gcc2}}\label{sec:proof}
%% ====================================================================

\subsection{Reduction to monotonicity}

With $s_i:=\tfrac{1}{2}t_i^2$, $K:=[0,s_1]\times\cdots\times[0,s_n]$,
and $Z(\tau)$ as above, \eqref{eq:gcc2} is equivalent to
$\PP(Z(1)\in K)\geqslant\PP(Z(0)\in K)$.  At $\tau=0$,
$C(0)=\diag(C_{11},C_{22})$ is block-diagonal, so the two blocks are
independent and the RHS factorizes.  At $\tau=1$ we recover $C$.
Hence \Cref{thm:gcc2} follows from:

\begin{proposition}\label{prop:key}
  $\tau\mapsto\int_K f(x,\tau)\,\dd x$ is non-decreasing on $[0,1]$.
\end{proposition}

\subsection{Laplace-transform computation of $\partial_\tau f$}

\Cref{prop:key} amounts to $\int_K\partial_\tau f(x,\tau)\,\dd x\geqslant 0$
for each $\tau\in(0,1)$ (\Cref{lem:diff}(ii) with $\lambda=0$ justifies
$\partial_\tau\int_K f=\int_K\partial_\tau f$).  Strategy:
\begin{enumerate}[label=(\alph*)]
  \item compute the Laplace transform of $\partial_\tau f$ and match
    it with that of a sum of derivatives of $h_\tau$;
  \item use Laplace injectivity (\Cref{thm:laplace_inj}) to get an
    a.e.\ pointwise identity;
  \item integrate this identity over $K$.
\end{enumerate}

\paragraph{Step (a).}  Since $Z_i(\tau)=\tfrac12 X_i(\tau)^2$, applying
\Cref{lem:Lap} to $X(\tau)\sim\caln(0,C(\tau))$ with $\lambda_i/2$ in
place of $\lambda_i$,
\[
  \int_{[0,\infty)^n}e^{-\sum_i\lambda_ix_i}f(x,\tau)\,\dd x
  =\EE e^{-\sum_i\lambda_iZ_i(\tau)}
  =|I_n+\Lambda C(\tau)|^{-1/2}.
\]
Differentiating in $\tau$, with \Cref{lem:diff}(ii) (taking $K=[0,\infty)^n$)
justifying the interchange:
\begin{equation}\label{eq:lt_deriv}
  \int_{[0,\infty)^n}e^{-\sum_i\lambda_ix_i}
  \partial_\tau f(x,\tau)\,\dd x
  =\partial_\tau|I_n+\Lambda C(\tau)|^{-1/2}.
\end{equation}
Expand via \Cref{lem:det}(i) applied to $A=\Lambda C(\tau)$:
\begin{equation}\label{eq:det_expansion}
  |I_n+\Lambda C(\tau)|
  =1+\sum_{\varnothing\ne J\subseteq[n]}|(\Lambda C(\tau))_J|
  =1+\sum_{\varnothing\ne J\subseteq[n]}|C(\tau)_J|\prod_{j\in J}\lambda_j,
\end{equation}
using $|(\Lambda C(\tau))_J|=|\Lambda_J||C(\tau)_J|=\prod_{j\in J}\lambda_j\cdot|C(\tau)_J|$
since $\Lambda$ is diagonal.

\paragraph{The block structure of $|C(\tau)_J|$.}  For $J\subseteq[n]$, set
$J_1:=[n_1]\cap J$ and $J_2:=J\setminus[n_1]$.  If $J_1=\varnothing$
or $J_2=\varnothing$, then $C(\tau)_J=C_J$ is $\tau$-independent,
since the off-diagonal blocks of $C(\tau)$ (where the $\tau$ appears)
are not sampled.  Otherwise, writing $C(\tau)_J$ in block form with
blocks indexed by $J_1,J_2$,
\[
  C(\tau)_J
  =\begin{pmatrix}C_{J_1}&\tau C_{J_1J_2}\\
    \tau C_{J_2J_1}&C_{J_2}\end{pmatrix},
\]
and applying \Cref{lem:det}(ii),
\begin{equation}\label{eq:CtauJ}
  |C(\tau)_J|
  =|C_{J_1}||C_{J_2}|\prod_{i=1}^{|J_1|}(1-\tau^2\mu_{J_1,J_2}(i)),
\end{equation}
where the $\mu_{J_1,J_2}(i)\in[0,1]$ are the eigenvalues of
$C_{J_1}^{-1/2}C_{J_1J_2}C_{J_2}^{-1}C_{J_2J_1}C_{J_1}^{-1/2}$
(bounded by \Cref{lem:det}(ii) since $C_J\succ 0$).
Differentiating the product in $\tau$,
\[
  \partial_\tau|C(\tau)_J|
  =-|C_{J_1}||C_{J_2}|\sum_{i=1}^{|J_1|}2\tau\mu_{J_1,J_2}(i)
  \prod_{l\ne i}(1-\tau^2\mu_{J_1,J_2}(l))
  \leqslant 0\quad\text{on }[0,1].
\]
Define
\begin{equation}\label{eq:aJ}
  a_J(\tau):=-\partial_\tau|C(\tau)_J|\geqslant 0
  \quad(\tau\in[0,1]),
\end{equation}
with $a_J:=0$ if $J_1=\varnothing$ or $J_2=\varnothing$.  Then
\begin{equation}\label{eq:partial_det}
  \partial_\tau|I_n+\Lambda C(\tau)|
  =-\sum_{\varnothing\ne J\subseteq[n]}a_J(\tau)\prod_{j\in J}\lambda_j.
\end{equation}
Apply the chain rule
$\partial_\tau(u^{-1/2})=-\tfrac12 u^{-3/2}\,\partial_\tau u$ to
$u=|I_n+\Lambda C(\tau)|$: \eqref{eq:lt_deriv} becomes
\[
  \int_{[0,\infty)^n}e^{-\sum_i\lambda_ix_i}
    \partial_\tau f(x,\tau)\,\dd x
  =-\tfrac12|I_n+\Lambda C(\tau)|^{-3/2}
    \partial_\tau|I_n+\Lambda C(\tau)|,
\]
and substituting \eqref{eq:partial_det} (the two minus signs cancel):
\begin{equation}\label{eq:lt_of_ftau}
  \int_{[0,\infty)^n}e^{-\sum_i\lambda_ix_i}
  \partial_\tau f(x,\tau)\,\dd x
  =\sum_{\varnothing\ne J\subseteq[n]}
  \tfrac{1}{2}a_J(\tau)|I_n+\Lambda C(\tau)|^{-3/2}
  \prod_{j\in J}\lambda_j.
\end{equation}

\subsection{Identification via $h_\tau:=h_{3,C(\tau)}$}

By \Cref{lem:LaplhkC} with $k=3$ and \Cref{lem:proph}(iii) (valid
since $\alpha_i=\tfrac{3}{2}>1$),
\[
  \int e^{-\sum_i\lambda_ix_i}
    \tfrac{\partial^{|J|}}{\partial x_J}h_\tau\,\dd x
  =\Bigl(\prod_{j\in J}\lambda_j\Bigr)|I_n+\Lambda C(\tau)|^{-3/2}.
\]
Comparing with \eqref{eq:lt_of_ftau}: $\partial_\tau f(\cdot,\tau)$
and $\sum_J\tfrac12 a_J(\tau)\,\partial^{|J|}h_\tau/\partial x_J$ are
$L_1$ with the same Laplace transform, so by \Cref{thm:laplace_inj},
\begin{equation}\label{eq:key_identity}
  \partial_\tau f(x,\tau)
  =\sum_{\varnothing\ne J\subseteq[n]}
    \tfrac{1}{2}a_J(\tau)
    \tfrac{\partial^{|J|}}{\partial x_J}h_\tau(x)
  \quad\text{a.e.\ on }(0,\infty)^n.
\end{equation}

\subsection{Integration over $K$}

Integrate \eqref{eq:key_identity} over $K=\prod_j[0,s_j]$.
For $J^c:=[n]\setminus J$ and $(s_J,x_{J^c})_i=s_i$ $(i\in J)$,
$=x_i$ $(i\in J^c)$, the fundamental theorem of calculus in each
$x_j$, $j\in J$, gives
\begin{equation}\label{eq:boundary_eval}
  \int_K\tfrac{\partial^{|J|}}{\partial x_J}h_\tau(x)\,\dd x
  =\int_{\prod_{j\in J^c}[0,s_j]}h_\tau(s_J,x_{J^c})\,\dd x_{J^c}
  \geqslant 0,
\end{equation}
since the lower-endpoint contributions vanish by
\Cref{lem:proph}(ii) ($\alpha_j=\tfrac{3}{2}>1$), and
$h_\tau\geqslant 0$.  Since also $a_J\geqslant 0$,
\[
  \int_K\partial_\tau f(x,\tau)\,\dd x
  =\sum_{\varnothing\ne J\subseteq[n]}\tfrac{1}{2}a_J(\tau)
  \int_{\prod_{j\in J^c}[0,s_j]}h_\tau(s_J,x_{J^c})\,\dd x_{J^c}
  \geqslant 0,
\]
which proves \Cref{prop:key} and hence \Cref{thm:gcc2}.  \qed


%% ====================================================================
\section{Two corollaries}\label{sec:corollaries}
%% ====================================================================

\begin{corollary}[\v{S}id\'ak]\label{cor:sidak}
  For $X\sim\caln(0,C)$ and $s_i\geqslant 0$,
  \[
    \PP(|X_i|\leqslant s_i,\;i=1,\ldots,n)
    \geqslant\prod_{i=1}^n\PP(|X_i|\leqslant s_i).
  \]
\end{corollary}

\begin{proof}
  Induct on $n$.  Case $n=1$ is trivial.  For the step, take
  $K_1:=\{|x_1|\leqslant s_1\}$ and
  $L_1:=\{|x_i|\leqslant s_i,2\leqslant i\leqslant n\}$ (both closed symmetric
  convex).  \Cref{thm:gcc2} (applied with $n_1=1$, $n_2=n-1$) gives
  $\PP(X\in K_1\cap L_1)\geqslant\PP(X\in K_1)\PP(X\in L_1)$;
  apply induction to $\PP(X\in L_1)$.
\end{proof}

\begin{corollary}[Decoupling]\label{cor:decouple}
  For $X\sim\caln(0,C)$ and closed symmetric convex $K,L\subseteq\R^n$,
  \[
    \PP(X\in K\cap L)\geqslant\PP(X\in K)\,\PP(X\in L).
  \]
\end{corollary}

\begin{proof}
  Restatement of \Cref{thm:gcc_geom} for this $K,L$.
\end{proof}

\end{document}
